\SourceLeaf{247}
\Article{235}
If the form
\[
AXX+2BXY+CYY\ldots F
\]
passes into the product of two forms
\[
axx+2bxy+cyy\ldots f,\quad\text{and}\quad a'x'x'+2b'x'y'+c'y'y'\ldots f'
\]
by a substitution of the kind
\[
\begin{aligned}
X&=pxx'+p'xy'+p''yx'+p'''yy'\\
Y&=qxx'+q'xy'+q''yx'+q'''yy'
\end{aligned}
\]
(which, for brevity, we will always express below as follows: If $F$ passes into $ff'$ by the substitution $p,p',p'',p''';\ q,q',q'',q'''$\HistoricalNote{*)}{In this notation, then, careful attention must be paid both to the order of the coefficients $p,p'$, etc., and to that of the forms $f,f'$. It is easy to see that if the order of the forms $f,f'$ is reversed, so that the first becomes the second, the coefficients $p',q'$ must be exchanged with $p'',q''$, while each of the others stays in its place.}), we will simply say that the form $F$ is \textit{transformable} into $ff'$; if, in addition, this transformation is such that the six numbers
\[
\begin{gathered}pq'-qp',\quad pq''-qp'',\quad pq'''-qp''',\\p'q''-q'p'',\quad p'q'''-q'p''',\quad p''q'''-q''p'''\end{gathered}
\]
have no common divisor, we will call the form $F$ \textit{composed} of the forms $f,f'$.

We begin this investigation with the most general assumption, that the form $F$ passes into $ff'$ by the substitution $p,p',p'',p''';\ q,q',q'',q'''$, and work out what follows from it. This assumption is clearly exactly equivalent to the following nine equations (\textit{i.e.}, as soon as these equations hold, $F$ passes into $ff'$ by the stated substitution, and conversely):
\begin{align*}
App+2Bpq+Cqq&=aa'\quad\ldotp\,\ldotp\,\ldotp\tag*{[1]}\\
Ap'p'+2Bp'q'+Cq'q'&=ac'\quad\ldotp\,\ldotp\,\ldotp\tag*{[2]}\\
Ap''p''+2Bp''q''+Cq''q''&=ca'\quad\ldotp\,\ldotp\,\ldotp\tag*{[3]}\\
Ap'''p'''+2Bp'''q'''+Cq'''q'''&=cc'\quad\ldotp\,\ldotp\,\ldotp\tag*{[4]}\\
App'+B(pq'+qp')+Cqq'&=ab'\quad\ldotp\,\ldotp\,\ldotp\tag*{[5]}\\
App''+B(pq''+qp'')+Cqq''&=ba'\quad\ldotp\,\ldotp\,\ldotp\tag*{[6]}\\
Ap'p'''+B(p'q'''+q'p''')+Cq'q'''&=bc'\quad\ldotp\,\ldotp\,\ldotp\tag*{[7]}\\
Ap''p'''+B(p''q'''+q''p''')+Cq''q'''&=cb'\quad\ldotp\,\ldotp\,\ldotp\tag*{[8]}\\
\begin{aligned}A(pp'''+p'p'')&+B(pq'''+qp'''+p'q''+q'p'')\\&+C(qq'''+q'q'')\end{aligned}&=2bb'\quad\ldotp\,\ldotp\,\ldotp\tag*{[9]}
\end{align*}

\SourceLeaf{248}
Let the determinants of the forms $F,f,f'$ be $D,d,d'$, respectively; let the greatest common divisors of the numbers $A,2B,C;\ a,2b,c;\ a',2b',c'$ be $M,m,m'$, respectively (all taken positively). Further, determine six integers $\mathfrak{A},\mathfrak{B},\mathfrak{C},\mathfrak{A}',\mathfrak{B}',\mathfrak{C}'$ so that
\[
\mathfrak{A}a+2\mathfrak{B}b+\mathfrak{C}c=m,\qquad
\mathfrak{A}'a'+2\mathfrak{B}'b'+\mathfrak{C}'c'=m'
\]
Finally, denote the numbers
\[
\begin{gathered}pq'-qp',\quad pq''-qp'',\quad pq'''-qp''',\\p'q''-q'p'',\quad p'q'''-q'p''',\quad p''q'''-q''p'''\end{gathered}
\]
by $P,Q,R,S,T,U$, respectively, and let their greatest common divisor, taken positively, be $=k$. --- Now, setting
\begin{align*}
App'''+B(pq'''+qp''')+Cqq'''&=bb'+\Delta\quad\ldotp\,\ldotp\,\ldotp\,\ldotp\tag*{[10]}
\end{align*}
we obtain from equation 9
\begin{align*}
Ap'p''+B(p'q''+q'p'')+Cq'q''&=bb'-\Delta\quad\ldotp\,\ldotp\,\ldotp\,\ldotp\tag*{[11]}
\end{align*}

From these eleven equations $1\ldots11$ we derive the following new ones\HistoricalNote{*)}{These equations arise as follows: 12 from $5.5-1.2$; 13 from $5.9-1.7-2.6$; 14 from $10.11-6.7$; 15 from $5.8+5.8+10.10+11.11-1.4-2.3-6.7-6.7$; 16 from $8.9-3.7-4.6$; 17 from $8.8-3.4$. The other six are derived in the same way, if we exchange equations $2,5,7$ with equations $3,6,8$, respectively, and keep the remaining ones $1,4,9,10,11$ in the same positions throughout; thus 18 comes from $6.6-1.3$, etc.}:
\begin{align*}
DPP&=d'aa\quad\ldotp\,\ldotp\,\ldotp\,\ldotp\,\ldotp\,\ldotp\,\ldotp\,\ldotp\,\ldotp\,\ldotp\,\ldotp\,\ldotp\,\ldotp\tag*{[12]}\\
DP(R-S)&=2d'ab\quad\ldotp\,\ldotp\,\ldotp\,\ldotp\,\ldotp\,\ldotp\,\ldotp\,\ldotp\,\ldotp\,\ldotp\,\ldotp\,\ldotp\tag*{[13]}\\
DPU&=d'ac-(\Delta\Delta-dd')\quad\ldotp\,\ldotp\,\ldotp\,\ldotp\,\ldotp\,\ldotp\,\ldotp\tag*{[14]}\\
D(R-S)^2&=4d'bb+2(\Delta\Delta-dd')\quad\ldotp\,\ldotp\,\ldotp\,\ldotp\,\ldotp\tag*{[15]}\\
D(R-S)U&=2d'bc\quad\ldotp\,\ldotp\,\ldotp\,\ldotp\,\ldotp\,\ldotp\,\ldotp\,\ldotp\,\ldotp\,\ldotp\,\ldotp\,\ldotp\tag*{[16]}\\
DUU&=d'cc\quad\ldotp\,\ldotp\,\ldotp\,\ldotp\,\ldotp\,\ldotp\,\ldotp\,\ldotp\,\ldotp\,\ldotp\,\ldotp\,\ldotp\,\ldotp\tag*{[17]}\\
DQQ&=da'a'\quad\ldotp\,\ldotp\,\ldotp\,\ldotp\,\ldotp\,\ldotp\,\ldotp\,\ldotp\,\ldotp\,\ldotp\,\ldotp\,\ldotp\tag*{[18]}\\
DQ(R+S)&=2da'b'\quad\ldotp\,\ldotp\,\ldotp\,\ldotp\,\ldotp\,\ldotp\,\ldotp\,\ldotp\,\ldotp\,\ldotp\,\ldotp\tag*{[19]}\\
DQT&=da'c'-(\Delta\Delta-dd')\quad\ldotp\,\ldotp\,\ldotp\,\ldotp\,\ldotp\,\ldotp\tag*{[20]}\\
D(R+S)^2&=4db'b'+2(\Delta\Delta-dd')\quad\ldotp\,\ldotp\,\ldotp\,\ldotp\,\ldotp\tag*{[21]}\\
D(R+S)T&=2db'c'\quad\ldotp\,\ldotp\,\ldotp\,\ldotp\,\ldotp\,\ldotp\,\ldotp\,\ldotp\,\ldotp\,\ldotp\,\ldotp\,\ldotp\tag*{[22]}\\
DTT&=dc'c'\quad\ldotp\,\ldotp\,\ldotp\,\ldotp\,\ldotp\,\ldotp\,\ldotp\,\ldotp\,\ldotp\,\ldotp\,\ldotp\,\ldotp\,\ldotp\tag*{[23]}
\end{align*}

From these we in turn deduce the following two:
\SourceLeaf{249}
\[
\begin{aligned}
0&=2d'aa(\Delta\Delta-dd')\\
0&=(\Delta\Delta-dd')^2-2d'ac(\Delta\Delta-dd')
\end{aligned}
\]
namely, the first from $12.15-13.13$, the second from $14.14-12.17$; from this it is easy to see that necessarily $\Delta\Delta-dd'=0$, whether $a=0$ or not $=0$\HistoricalNote{*)}{This derivation of the equation $\Delta\Delta=dd'$ is enough for our present purpose; otherwise we could give a more elegant analysis, but one too long for this place, by directly deriving from equations $1\ldots11$ the equation $0=(\Delta\Delta-dd')^2$.}. So we will suppose that in equations $14,15,20,21$, $\Delta\Delta-dd'$ is removed from the right sides.

Now, setting
\[
\begin{aligned}
\mathfrak{A}P+\mathfrak{B}(R-S)+\mathfrak{C}U&=mn'\\
\mathfrak{A}'Q+\mathfrak{B}'(R+S)+\mathfrak{C}'T&=m'n
\end{aligned}
\]
(where we must carefully note that $n,n'$ can also be fractions, although $mn',m'n$ are necessarily integers), we readily deduce from equations $12\ldots17$
\[
Dmmn'n'=d'(\mathfrak{A}a+2\mathfrak{B}b+\mathfrak{C}c)^2=d'mm
\]
and similarly from equations $18\ldots23$
\[
Dm'm'nn=d(\mathfrak{A}'a'+2\mathfrak{B}'b'+\mathfrak{C}'c')^2=dm'm'
\]

Thus $d=Dnn,\ d'=Dn'n'$, and we obtain the \textsc{first conclusion}: \textit{The determinants of the forms $F,f,f'$ necessarily have the ratios of squares to one another;} and the \textsc{second}: \textit{$D$ always divides the numbers $dm'm',d'mm$.} So $D,d,d'$ clearly have the same sign, and no form can be transformable into the product $ff'$ if its determinant is greater than the greatest common divisor of the numbers $dm'm',d'mm$.

Multiply equations $12,13,14$ by $\mathfrak{A},\mathfrak{B},\mathfrak{C}$, respectively; likewise multiply equations $13,15,16$, and $14,16,17$, by the same numbers; add each three products, and divide the sum by $Dmn'$, replacing $d'$ by $Dn'n'$. We then get
\[
P=an',\quad R-S=2bn',\quad U=cn'
\]

Similarly, multiplying equations $18,19,20$, as well as $19,21,22$ and $20,22,23$, by $\mathfrak{A}',\mathfrak{B}',\mathfrak{C}'$, respectively, we obtain
\[
Q=a'n,\quad R+S=2b'n,\quad T=c'n
\]

\SourceLeaf{250}
This gives the \textsc{third conclusion}: \textit{The numbers $a,2b,c$ are proportional to the numbers $P,R-S,U$, and if their ratio to the latter is put as 1 to $n'$, then $n'$ is a square root of $\frac{d'}{D}$; similarly, the numbers $a',2b',c'$ have the same ratio to $Q,R+S,T$, and if this is put as 1 to $n$, then $n$ is a square root of $\frac{d}{D}$.}

The quantities $n,n'$, however, may be either positive or negative square roots of $\frac{d}{D},\frac{d'}{D}$. This gives us a distinction that may seem fruitless at first, but whose use will become clear enough below. Namely, in the transformation of the form $F$ into $ff'$, we will say that the form $f$ is taken \textit{directly} when $n$ is positive, and \textit{inversely} when $n$ is negative; similarly, $f'$ is taken directly or inversely according as $n'$ is positive or negative. If we also require that $k$ be $=1$, the form $F$ will be said to be composed \textit{either} directly of both forms $f,f'$, \textit{or} inversely of both, \textit{or} directly of $f$ and inversely of $f'$, \textit{or} inversely of $f$ and directly of $f'$, according as $n,n'$ are both positive, both negative, the first positive and the second negative, or the first negative and the second positive. Anyone will readily understand that these relations do not depend on the order in which the forms $f,f'$ are placed (see the first note to this article).

We further observe that the greatest common divisor of the numbers $P,\allowbreak Q,\allowbreak R,\allowbreak S,\allowbreak T,\allowbreak U$, namely $k$, divides the numbers $mn',m'n$ (as is clear from the values established above), and so the square $kk$ divides $mmn'n',m'm'nn$, and $Dkk$ divides $d'mm,dm'm'$. Conversely, every common divisor of $mn',m'n$ also divides $k$. For let $e$ be such a divisor; it clearly also divides the numbers $an',2bn',cn',a'n,2b'n,c'n$, \textit{i.e.}, the numbers $P,R-S,U,Q,R+S,T$, and therefore also $2R$ and $2S$. Now if $\frac{2R}{e}$ were odd, $\frac{2S}{e}$ would have to be odd too (since their sum and difference are even), and so their product would also be odd. But this product is $=\frac{4}{ee}(b'b'nn-bbn'n')=\frac{4}{ee}(d'nn+a'c'nn-dn'n'-acn'n')=\frac{4}{ee}(a'c'nn-acn'n')$ and is therefore even, since $e$ divides $a'n,c'n,an',cn'$. So $\frac{2R}{e}$ is necessarily even, and therefore $R$, as well as $S$, is divisible by $e$. Since $e$ thus divides all six numbers $P,Q,R,S,T,U$, it also divides their greatest common divisor $k$. \textit{Q. E. D.} --- We conclude that $k$ is the greatest common divisor of the numbers $mn',m'n$; from this it is easy to see that \textit{$Dkk$ is the greatest common divisor of the numbers $dm'm',d'mm$.} This is the \textsc{fourth conclusion}. So whenever $F$ is composed of $f$ and $f'$, $D$ is clearly the greatest common divisor of the numbers $dm'n',d'mm$, and conversely; this property could also have been adopted as the definition of a composed form. Thus a form composed of the
\SourceLeaf{251}
forms $f,f'$ has the greatest possible determinant among all forms transformable into the product $ff'$.

Before we can proceed further, we must first determine more precisely the value of $\Delta$, which we have shown to be $=\surd dd'=\surd DDnnn'n'$, but whose \textit{sign} has not yet been determined by this. For this purpose, from the fundamental equations 1---11 we derive $DPQ=\Delta aa'$ (this equation comes from $5.6-1.11$), and so $Daa'nn'=\Delta aa'$, from which, unless one of the numbers $a,a'$ is $=0$, we get $\Delta=Dnn'$. But in exactly the same way, eight other equations can be derived from the fundamental equations, in which $Dnn'$ on the left and $\Delta$ on the right are multiplied by $2ab',ac',2ba',4bb',2bc',ca',2cb',cc'$\HistoricalNote{*)}{For brevity we must leave out the analysis, which readers can easily find for themselves.}. From this we easily conclude, since neither all of $a,2b,c$ nor all of $a',2b',c'$ can be $=0$, that in all cases $\Delta=Dnn'$, and therefore $\Delta$ has the same sign as $D,d,d'$, or the opposite sign, according as $n,n'$ have the same or different signs.

We further observe that the numbers $aa',2ab',ac',2ba',4bb',2bc',ca',2cb',cc',2bb'+2\Delta,2bb'-2\Delta$ are all divisible by $mm'$. This is evident for the first nine; for the other two, it can be proved in the same way as we showed above that $R$ and $S$ are divisible by $e$. Namely, $4bb'+4\Delta$ and $4bb'-4\Delta$ are clearly divisible by $mm'$ (since $4\Delta=\surd16dd'$, and $4d$ is divisible by $mm$, $4d'$ by $m'm'$, and therefore $16dd'$ by $mmm'm'$ and $4\Delta$ by $mm'$), and the difference of the quotients is even; their product is easily shown to be even, so each quotient is even, and $2bb'+2\Delta,2bb'-2\Delta$ are divisible by $mm'$.

The following six equations are now easily deduced from the eleven fundamental equations:
\[
\begin{aligned}
APP&=aa'q'q'-2ab'qq'+ac'qq\\
AQQ&=aa'q''q''-2ba'qq''+ca'qq\\
ARR&=aa'q'''q'''-2(bb'+\Delta)qq'''+cc'qq\\
ASS&=ac'q''q''-2(bb'-\Delta)q'q''+ca'q'q'\\
ATT&=ac'q'''q'''-2bc'q'q'''+cc'q'q'\\
AUU&=ca'q'''q'''-2cb'q''q'''+cc'q''q''
\end{aligned}
\]

It follows that all of $APP,AQQ$, etc. are divisible by $mm'$, and from this, since $kk$ is the greatest common divisor of the numbers $PP,QQ,$
\SourceLeaf{252}
$RR$, etc., it is easy to deduce that $Akk$ is also divisible by $mm'$. If we substitute for $a,2b,c$, $a',2b',c'$ their values $\frac{P}{n'}$, etc., or $\frac{1}{n'}(pq'-qp')$, etc., these become six other equations, in which the right sides are products of the quantity $\frac{1}{nn'}(q'q''-qq''')$ with $PP,QQ,RR$, etc. We leave this very easy calculation to the readers. It follows (since not all of $PP,QQ$, etc. can be $=0$) that $Ann'=q'q''-qq'''$.

Similarly, six other equations are derived from the fundamental equations, differing from the previous ones only in that $A$ is everywhere replaced by $C$, and $q,q',q'',q'''$ by $p,p',p'',p'''$, respectively. We do not write them down, for brevity. In the same way it follows that $Ckk$ is divisible by $mm'$ and that $Cnn'=p'p''-pp'''$.

Finally, the same source gives these six equations:
\[
\begin{aligned}
BPP&=-aa'p'q'+ab'(pq'+qp')-ac'pq\\
BQQ&=-aa'p''q''+ba'(pq''+qp'')-ca'pq\\
BRR&=-aa'p'''q'''+(bb'+\Delta)(pq'''+qp''')-cc'pq\\
BSS&=-ac'p''q''+(bb'-\Delta)(p'q''+q'p'')-ca'p'q'\\
BTT&=-ac'p'''q'''+bc'(p'q'''+q'p''')-cc'p'q'\\
BUU&=-ca'p'''q'''+cb'(p''q'''+q''p''')-cc'p''q''
\end{aligned}
\]
from which we conclude, just as before, that $2Bkk$ is divisible by $mm'$ and that $2Bnn'=pq'''+qp'''-p'q''-q'p''$.

Since $Akk,2Bkk,Ckk$ are therefore divisible by $mm'$, we readily see that $Mkk$ must also be divisible by $mm'$. But the fundamental equations show that $M$ divides $aa',\allowbreak 2ab',\allowbreak ac',\allowbreak 2ba',\allowbreak 4bb',\allowbreak 2bc',\allowbreak ca',\allowbreak 2cb',\allowbreak cc'$, and therefore also $am',2bm',cm'$ (the greatest common divisors of the first, middle and last three, respectively); finally, it also divides $mm'$, their greatest common divisor. So in the case where the form $F$ is composed of the forms $f,f'$, or $k=1$, we necessarily have $M=mm'$. This is the \textsc{fifth conclusion}.

If the greatest common divisor of the numbers $A,B,C$ is $\mathfrak{M}$, it is either $=M$ (when the form $F$ is properly primitive or derived from a properly primitive form) or $=\frac{1}{2}M$ (when $F$ is improperly primitive or derived from an improperly primitive form); similarly, if we denote the greatest common divisors of the numbers $a,b,c;\ a',b',c'$ by $\mathfrak{m},\mathfrak{m}'$, respectively, then $\mathfrak{m}$ is either $=m$ or $\frac{1}{2}m$, and $\mathfrak{m}'$ either $=m'$ or $=\frac{1}{2}m'$. Now clearly $\mathfrak{m}\mathfrak{m}$ divides $d'$, $\mathfrak{m}'\mathfrak{m}'$ divides $d'$, and therefore $\mathfrak{m}\mathfrak{m}\mathfrak{m}'\mathfrak{m}'$ divides $dd'$ or $\Delta\Delta$, and $\mathfrak{m}\mathfrak{m}'$ divides $\Delta$.
\SourceLeaf{253}
From the last six equations for $BPP$, etc. it then follows that $\mathfrak{m}\mathfrak{m}'$ divides $Bkk$ and so (since it also divides $Akk,Ckk$) also $\mathfrak{M}kk$. Thus whenever $F$ is composed of $f,f'$, $\mathfrak{m}\mathfrak{m}'$ divides $\mathfrak{M}$. So when, in this case, both $f,f'$ are properly primitive or derived from properly primitive forms, or $\mathfrak{m}\mathfrak{m}'=mm'=M$, we have $\mathfrak{M}=M$, or $F$ is a similar form. But when, under the same assumption, either both $f,f'$ or at least one of them is improperly primitive or derived from an improperly primitive form, for \textit{example} the form $f$, the fundamental equations show that $aa',2ab',ac',ba',2bb',bc',ca',2cb',cc'$ are divisible by $\mathfrak{M}$, and therefore also $am',bm',cm'$, and consequently $\mathfrak{m}m'=\frac{1}{2}mm'=\frac{1}{2}M$ too. So in this case we necessarily have $\mathfrak{M}=\frac{1}{2}M$, or the form $F$ is also improperly primitive or derived from an improperly primitive form. These give the \textsc{sixth conclusion}.

Finally, we observe that if the nine equations
\[
\begin{gathered}
an'=P,\quad 2bn'=R-S,\quad cn'=U\\
a'n=Q,\quad 2b'n=R+S,\quad c'n=T\\
\begin{gathered}Ann'=q'q''-qq''',\\2Bnn'=pq'''+qp'''-p'q''-q'p'',\\Cnn'=p'p''-pp'''\end{gathered}
\end{gathered}
\]
(which we will denote by $\Omega$, since we will often need to return to them below) are \textit{assumed} to hold, even treating $n,n'$ as unknowns, but neither being $=0$, substitution readily confirms that the fundamental equations $1\ldots9$ are necessarily true as well, or that the form $(A,B,C)$ passes by the substitution $p,p',p'',p''';\ q,q',q'',q'''$ into the product of the forms $(a,b,c)(a',b',c')$; and also that
\[
bb-ac=nn(BB-AC),\quad b'b'-a'c'=n'n'(BB-AC)
\]
We leave to the readers' efforts the calculation, which would be too long to include here.
